\documentclass[conference]{IEEEtran}

\usepackage{iftex}
\ifPDFTeX
  \usepackage[utf8]{inputenc}
  \usepackage[T1]{fontenc}
\fi

\usepackage{amsmath,bm,mathtools}
\usepackage{newtxtext,newtxmath}
\usepackage{textcomp}
\usepackage{enumitem}
\usepackage{booktabs}
\usepackage{siunitx}
\usepackage{float}
\usepackage[breaklinks=true,hidelinks]{hyperref}
\usepackage[nameinlink,capitalise,noabbrev]{cleveref}
\usepackage[final]{microtype}

\interdisplaylinepenalty=2500
\allowdisplaybreaks

\DeclareMathOperator{\clamp}{clamp}
\DeclareMathOperator{\atanTwo}{atan2}

\newcommand{\vct}[1]{\bm{#1}}
\newcommand{\R}{\mathbb{R}}
\newcommand{\dd}{\,\mathrm{d}}

\title{Adaptive Vertical PII with Mahony Attitude and Frequency-Based Scheduling\\
(Implementation Note)}

\author{%
  \IEEEauthorblockN{Mikhail S.\ Grushinskiy}
  \IEEEauthorblockA{Bareboat Necessities Project\\
  \texttt{mgrouch@users.github.com}}
}

\begin{document}
\maketitle

\begin{abstract}
This note documents the method implemented by
\texttt{AdaptiveVerticalPIIMahony.h},
\texttt{AdaptiveVerticalPII.h}, and
\texttt{VerticalPIIObserver.h}, together with the simulation wrapper
\texttt{w3d\_sim\_adaptive\_pii\_mahony.cpp}.
The pipeline estimates attitude with a Mahony complementary filter,
rotates body-frame specific force into a \(Z\)-up world frame,
subtracts gravity to obtain vertical inertial acceleration,
and drives an adaptive vertical displacement observer whose gains are
scheduled from a frequency proxy and an acceleration-variance estimate.

The implementation is layered.
First, a repeated-real-pole vertical PII observer estimates
displacement and velocity from world-frame vertical inertial acceleration.
Second, a Mahony wrapper supplies that vertical acceleration from raw IMU data.
Third, an adaptive wrapper schedules observer parameters from a policy-selected
frequency tracker and a running acceleration sigma estimate.
The emphasis here is on the actual implemented equations and signal flow.
\end{abstract}

\begin{IEEEkeywords}
heave estimation, vertical observer, PII observer, Mahony AHRS,
adaptive scheduling, marine inertial sensing
\end{IEEEkeywords}

% ---------------------------------------------------------------
\section{Signal Flow and Layering}

The implementation is organized as a small stack:
\begin{enumerate}[leftmargin=1.5em]
  \item \texttt{VerticalPIIObserver}: core vertical-motion observer,
  \item \texttt{AdaptiveVerticalPII}: wrapper adding sigma estimation,
        frequency-proxy scheduling, and tracker-policy abstraction,
  \item \texttt{AdaptiveVerticalPIIMahony}: wrapper that uses Mahony AHRS
        to convert raw IMU data into vertical inertial acceleration.
\end{enumerate}

The simulation entry point instantiates
\texttt{AdaptiveVerticalPIIMahony<float,true,TrackerType::ARANOVSKIY>}
in the supplied driver, but the adaptive wrapper itself is generic and can
use several tracker backends through \texttt{FrequencyTrackerPolicy}.

At sample time \(k\) with period \(\Delta t_k\), the implemented cascade is
\begin{align}
\vct{\omega}_{b,k} &\xrightarrow{\text{Mahony}} \hat{q}_{wb,k}, \\
\vct{f}_{b,k} &\xrightarrow{\hat{q}_{bw,k}} \hat{\vct{f}}_{w,k}, \\
\hat{a}_{z,k} &= \hat{f}_{w,z,k} - g, \\
\hat{a}_{z,k} &\xrightarrow{\text{Adaptive Vertical PII}}
(\hat{p}_k,\hat{v}_k,\hat{S}_k,\hat{b}_k).
\end{align}
Here:
\begin{itemize}[leftmargin=1.4em]
  \item \(\hat{q}_{wb}\) is the Mahony world-to-body quaternion,
  \item \(\hat{q}_{bw}=\hat{q}_{wb}^{\ast}\) is its conjugate,
  \item \(\hat{p}\) is vertical displacement,
  \item \(\hat{v}\) is vertical velocity,
  \item \(\hat{S}\) is the integral of displacement,
  \item \(\hat{b}\) is the optional slow acceleration-bias estimate.
\end{itemize}

The rest of this note follows the code structure requested here:
first the vertical PII observer, then the Mahony wrapper,
then the adaptive scheduling layer.

% ---------------------------------------------------------------
\section{Vertical PII Observer}\label{sec:pii}

\subsection{Purpose and States}

\texttt{VerticalPIIObserver} estimates vertical displacement and velocity from
an input that is already the gravity-compensated world-frame vertical inertial
acceleration.

Its main states are:
\begin{align}
a_f &: \text{filtered vertical acceleration},\\
v   &: \text{vertical velocity},\\
p   &: \text{vertical displacement},\\
S   &: \text{integral of displacement}.
\end{align}

When the optional bias channel is enabled, it also keeps
\begin{align}
d &: \text{very slow displacement trend},\\
b &: \text{estimated acceleration bias}.
\end{align}

\subsection{Continuous-Time Model}

The observer equations documented in the header are
\begin{align}
\dot{a}_f &= \frac{a_{\mathrm{meas}} - a_f}{\tau_a}, \label{eq:afdot}\\
\dot{S}   &= p, \label{eq:Sdot}\\
\dot{p}   &= v, \label{eq:pdot}\\
\dot{v}   &= a_f - b - k_v v - k_p p - k_s S. \label{eq:vdot}
\end{align}

When \texttt{WithBias=true}, the very-slow bias-trend channel is
\begin{align}
\dot{d} &= \frac{p-d}{\tau_d}, \label{eq:ddot}\\
\dot{b} &= k_b d - \lambda_b b. \label{eq:bdot}
\end{align}

So the observer is not just a double integrator. It is a damped
third-order restoring system driven by filtered acceleration, with an
optional very-slow leak-stabilized bias estimator.

\subsection{Repeated-Real-Pole Parameterization}

The core PII restoring loop is parameterized by a repeated real pole at
\(-r\):
\begin{equation}
(s+r)^3 = s^3 + 3rs^2 + 3r^2 s + r^3.
\end{equation}

Matching coefficients yields
\begin{equation}
k_v = 3r,\qquad
k_p = 3r^2,\qquad
k_s = r^3.
\end{equation}

This matters because the observer can be tuned through one physically
meaningful bandwidth-like parameter \(r\), while preserving an overdamped,
all-real-pole structure. The implementation derives the active gains
\((k_v,k_p,k_s)\) from the currently active \(r\) every time the active
motion-pole rate changes.

\subsection{Discrete Update Used by the Code}

The implementation uses the embedded-friendly one-pole coefficient
\begin{equation}
\alpha(dt,\tau) = \frac{dt}{\tau + dt},
\end{equation}
rather than an exponential form.

The filtered acceleration update is
\begin{equation}
a_f^{+} = a_f + \alpha(dt,\tau_a)(a_{\mathrm{meas}}-a_f).
\end{equation}

If the bias channel is enabled, the slow trend update is
\begin{equation}
d^{+} = d + \alpha(dt,\tau_d)(p-d),
\end{equation}
and the bias state is integrated as
\begin{equation}
b^{+} = b + dt\,(k_b d^{+} - \lambda_b b).
\end{equation}

The vertical PII core then forms
\begin{equation}
\dot{v} = a_f^{+} - b^{+} - k_v v - k_p p - k_s S,
\end{equation}
and applies a semi-implicit update:
\begin{align}
v^{+} &= v + dt\,\dot{v}, \\
p^{+} &= p + dt\,v^{+}, \\
S^{+} &= S + dt\,p^{+}.
\end{align}

That update order is important: the new velocity is used immediately in the
position update, and the new position is used immediately in the integral
update.

\subsection{State and Safety Clamps}

The implementation supports optional symmetric clamps on
\begin{equation}
a_f,\quad v,\quad p,\quad S,\quad d,\quad b.
\end{equation}
A non-positive limit disables the corresponding clamp. These are practical
safety bounds rather than part of the nominal observer model.

\begin{table}[t]
  \centering
  \caption{Main states of \texttt{VerticalPIIObserver}}
  \label{tab:pii_states}
  \small
  \begin{tabular}{@{}lll@{}}
    \toprule
    Quantity & Code field & Meaning \\
    \midrule
    \(a_f\) & \texttt{a\_f\_} & filtered vertical acceleration \\
    \(v\)   & \texttt{v\_}    & vertical velocity \\
    \(p\)   & \texttt{p\_}    & vertical displacement \\
    \(S\)   & \texttt{S\_}    & integral of displacement \\
    \(d\)   & \texttt{bias\_.d} & slow displacement trend \\
    \(b\)   & \texttt{bias\_.b} & estimated acceleration bias \\
    \(r\)   & \texttt{active\_r\_} & active repeated pole rate \\
    \bottomrule
  \end{tabular}
\end{table}

% ---------------------------------------------------------------
\section{Mahony Attitude Wrapper}\label{sec:mahony}

\subsection{World-Frame Convention}

The supplied Mahony implementation is not NED. In this code base it is used
as a \(Z\)-up auxiliary/world frame. At identity quaternion, the expected
accelerometer reading is approximately \((0,0,+g)\), so world up is \(+Z\).

Accordingly, the wrapper computes
\begin{equation}
a_{z,\mathrm{world,up}}
=
\bigl(R_{bw} a_b\bigr)_z - g,
\end{equation}
where:
\begin{itemize}[leftmargin=1.4em]
  \item \(a_b\) is body-frame specific force in \si{m/s^2},
  \item \(R_{bw}\) is the body-to-world rotation from Mahony,
  \item \(g\) is the configured gravity magnitude.
\end{itemize}

This vertical inertial acceleration is exactly the signal fed into the
adaptive vertical observer.

\subsection{Mahony Innovation Structure}

The Mahony state stores the world-to-body quaternion
\begin{equation}
\hat{q}_{wb} = [q_0\;q_1\;q_2\;q_3]^\top.
\end{equation}

In the IMU-only branch, the accelerometer is normalized and the predicted
body-frame direction of world \(+Z\) is extracted from the current quaternion
as
\begin{equation}
\hat{\vct{g}}_b =
\begin{bmatrix}
2(q_1 q_3 - q_0 q_2)\\
2(q_0 q_1 + q_2 q_3)\\
q_0^2 - q_1^2 - q_2^2 + q_3^2
\end{bmatrix}.
\end{equation}

If \(\hat{\vct{a}}_b\) is the normalized measured accelerometer direction, the
gravity innovation is
\begin{equation}
\vct{e} = \hat{\vct{a}}_b \times \hat{\vct{g}}_b.
\end{equation}

With magnetometer aiding, the same gravity correction is supplemented by a
magnetic innovation term inside the Mahony update routine.

\subsection{Mahony Kinematics}

The continuous-time idea is the standard Mahony structure:
\begin{equation}
\vct{\omega}_c = \vct{\omega}_m + \vct{i}_{\omega} + 2K_p\,\vct{e},
\end{equation}
with integral feedback
\begin{equation}
\dot{\vct{i}}_{\omega} = 2K_i\,\vct{e},
\end{equation}
and quaternion kinematics
\begin{equation}
\dot{\hat{q}}_{wb}
=
\frac12\,\hat{q}_{wb}\otimes [0\;\vct{\omega}_c].
\end{equation}

In code, this is implemented by the supplied Mahony functions using
per-sample quaternion updates and renormalization.

\subsection{Body-to-World Rotation}

The wrapper needs world-frame acceleration, so it conjugates the quaternion:
\begin{equation}
\hat{q}_{bw} = \hat{q}_{wb}^{\ast}.
\end{equation}

The code then rotates a body-frame vector into the Mahony world frame using
the efficient quaternion-vector formula
\begin{align}
\vct{t} &= 2\,\vct{q}_v \times \vct{f}_b,\\
\hat{\vct{f}}_w &= \vct{f}_b + q_w\vct{t} + \vct{q}_v \times \vct{t},
\end{align}
where \(\hat{q}_{bw}=[q_w\;\vct{q}_v]\).

The resulting world-frame specific force is
\begin{equation}
\hat{\vct{f}}_w = [\hat{f}_{w,x}\;\hat{f}_{w,y}\;\hat{f}_{w,z}]^\top,
\end{equation}
and the vertical inertial acceleration passed downstream is
\begin{equation}
\hat{a}_z = \hat{f}_{w,z} - g.
\end{equation}

At rest and level, \(\hat{f}_{w,z}\approx g\), so \(\hat{a}_z \approx 0\).

\subsection{Mahony Gain Adaptation}

The wrapper can also adapt Mahony gains according to how trustworthy the
accelerometer looks as a gravity cue. It computes:
\begin{align}
\varepsilon_{\mathrm{norm}} &=
\frac{\left|\|\vct{a}_b\| - g\right|}{g},\\
\varepsilon_{\mathrm{innov}} &=
\left\|\hat{\vct{a}}_b \times \hat{\vct{g}}_b\right\|,\\
\sigma_a &= \text{vertical acceleration sigma from the adaptive core}.
\end{align}

Those are normalized by configured reference scales and converted into trust
factors
\begin{equation}
\rho_n = \frac{1}{1+n^2},\qquad
\rho_e = \frac{1}{1+u^2},\qquad
\rho_s = \frac{1}{1+s^2},
\end{equation}
where \(n\), \(u\), and \(s\) are the normalized norm error, innovation,
and vertical-acceleration sigma.

The combined accelerometer trust is
\begin{equation}
\gamma = \clamp(\rho_n\,\rho_e\,\rho_s,\gamma_{\min},1).
\end{equation}

That trust interpolates between rough-water and calm-water Mahony gains:
\begin{align}
2K_p(\gamma)
&=
2K_{p,\mathrm{rough}} +
\gamma\bigl(2K_{p,\mathrm{calm}}-2K_{p,\mathrm{rough}}\bigr),\\
2K_i(\gamma)
&=
2K_{i,\mathrm{rough}} +
\gamma^2\bigl(2K_{i,\mathrm{calm}}-2K_{i,\mathrm{rough}}\bigr).
\end{align}

So the proportional term fades linearly with trust, while the integral term
collapses even faster in rough motion. The commanded gains are then smoothed
with the same one-pole coefficient form before being written back into the
Mahony state.

% ---------------------------------------------------------------
\section{Adaptive Scheduling Wrapper}\label{sec:adaptive}

\subsection{Role of \texttt{AdaptiveVerticalPII}}

\texttt{AdaptiveVerticalPII} wraps the core observer and adds:
\begin{itemize}[leftmargin=1.4em]
  \item a running estimate of vertical-acceleration mean, variance, and sigma,
  \item a policy-selected acceleration-frequency tracker,
  \item automatic or external scheduling calls into
        \texttt{VerticalPIIObserver::update\_adaptation()}.
\end{itemize}

It therefore sits between the raw vertical acceleration signal and the
observer’s scheduling interface.

\subsection{Acceleration Variance Estimate}

At IMU rate, the wrapper updates a slow mean and a shorter-time variance:
\begin{align}
\mu_a^{+} &= \mu_a + \alpha_\mu(a_z-\mu_a),\\
e_a &= a_z - \mu_a^{+},\\
\nu_a^{+} &= \nu_a + \alpha_\nu(e_a^2 - \nu_a),\\
\sigma_a^{+} &= \sqrt{\max(\nu_a^{+},0)}.
\end{align}

A configured floor is then enforced:
\begin{equation}
\sigma_a^{+} \leftarrow \max(\sigma_a^{+}, \sigma_{\min}).
\end{equation}

This \(\sigma_a\) is the roughness indicator used in observer scheduling and
also in Mahony trust adaptation.

\subsection{Scheduling Hook Inside \texttt{VerticalPIIObserver}}

Whenever adaptation is enabled and a scheduling update is triggered,
\texttt{VerticalPIIObserver::update\_adaptation()} receives:
\begin{equation}
(f_{\mathrm{disp}},\sigma_a,c,dt_{\mathrm{sched}}).
\end{equation}

If confidence is below the configured minimum, or if the inputs are invalid,
the active parameters are held.

Otherwise the observer first smooths the tracker inputs:
\begin{align}
f_f^{+} &= f_f + \alpha_{\mathrm{in}}(f_{\mathrm{disp}}-f_f),\\
\sigma_f^{+} &= \sigma_f + \alpha_{\mathrm{in}}(\sigma_a-\sigma_f).
\end{align}

Then it forms normalized scheduling ratios
\begin{equation}
\rho_f = \clamp\!\left(\frac{f_f}{f_{\mathrm{ref}}},0.25,4\right),
\qquad
\rho_\sigma = \clamp\!\left(\frac{\sigma_{\mathrm{ref}}}{\sigma_f},0.25,4\right).
\end{equation}

The commanded observer parameters are conservative power-law schedules:
\begin{align}
r_{\mathrm{cmd}} &=
r_0\,\rho_f^{\eta_r}\,\rho_\sigma^{\zeta_r},\\
\tau_{a,\mathrm{cmd}} &=
\tau_{a,0}\,\rho_f^{\eta_a}\,\rho_\sigma^{\zeta_a},\\
\tau_{d,\mathrm{cmd}} &=
\tau_{d,0}\,\rho_f^{\eta_d}\,\rho_\sigma^{\zeta_d},\\
k_{b,\mathrm{cmd}} &=
k_{b,0}\,\rho_f^{\eta_b}\,\rho_\sigma^{\zeta_b}.
\end{align}

Each commanded parameter is clamped to configured bounds, then the active
parameters move toward the commands via one-pole smoothing:
\begin{equation}
\theta^{+}
=
\theta + \alpha_p(\theta_{\mathrm{cmd}}-\theta),
\end{equation}
for each scheduled parameter \(\theta\).

\subsection{Interpretation of the Scheduling Signs}

The exponent signs encode the intended control effect:
\begin{itemize}[leftmargin=1.4em]
  \item higher dominant displacement frequency generally allows a somewhat
        faster observer core,
  \item larger acceleration sigma tends to reduce aggressiveness and slow bias
        learning,
  \item lower sigma permits more restoring authority.
\end{itemize}

So the scheduler is meant to adapt slowly to sea-state changes, not to chase
sample-by-sample motion.

% ---------------------------------------------------------------
\section{Frequency Proxy Used for Scheduling}\label{sec:freq_proxy}

\subsection{Generic Policy Interface}

\texttt{AdaptiveVerticalPII} does not hard-code one tracker algorithm.
Instead it uses a policy-selected backend
\begin{equation}
\texttt{AccelFreqTracker = TrackerPolicy<TT>},
\end{equation}
where \texttt{TT} is a template parameter such as PLL, Aranovskiy, KalmANF,
or ZeroCross.

Through the policy helpers, the wrapper queries:
\begin{itemize}[leftmargin=1.4em]
  \item main frequency estimate,
  \item raw frequency estimate,
  \item tracker confidence,
  \item lock status,
  \item optional coarse estimate.
\end{itemize}

\subsection{Automatic Scheduling from the Frequency Proxy}

When \texttt{auto\_schedule\_from\_accel\_freq} is enabled, the wrapper
accumulates time and periodically triggers an observer scheduling update.

The scheduled frequency is built by:
\begin{enumerate}[leftmargin=1.5em]
  \item preferring the tracker’s main frequency estimate,
  \item falling back to its raw frequency estimate if needed,
  \item blending toward a coarse estimate when one exists,
  \item blending more strongly toward the coarse estimate when the tracker is
        not locked.
\end{enumerate}

So the scheduling frequency is intentionally more conservative than just the
raw tracker output.

\subsection{Confidence Passed to the Scheduler}

The confidence passed to the observer scheduler begins with the tracker’s own
confidence. It is then floored upward when:
\begin{itemize}[leftmargin=1.4em]
  \item a coarse estimate exists,
  \item the tracker reports locked.
\end{itemize}

This lets the observer continue adapting moderately even when the internal
frequency proxy is only coarse-valid rather than fully locked.

\subsection{Simulation-Specific Note}

In the supplied simulation driver,
\texttt{FusionAdapterAdaptivePIIMahony} instantiates
\begin{equation}
\texttt{AdaptiveVerticalPIIMahony<float,true,TrackerType::ARANOVSKIY>}.
\end{equation}

So that file uses the Aranovskiy tracker as the internal acceleration-frequency
proxy, even though the adaptive wrapper itself remains generic.

% ---------------------------------------------------------------
\section{Simulation-Backed Tuning Example}\label{sec:tuning}

The simulation wrapper sets the following representative values for the core
observer:
\begin{align}
r &= 0.150,\\
\tau_a &= 0.68~\si{s},\\
\tau_d &= 49.0~\si{s},\\
k_b &= 2.5\times10^{-5},\\
\lambda_b &= 3.0\times10^{-3}.
\end{align}

Safety limits are
\begin{equation}
|a_f|\le 50,\quad |v|\le 50,\quad |p|\le 20,\quad |S|\le 200,\quad |d|\le 20,
\end{equation}
with
\begin{equation}
|b|\le 0.12.
\end{equation}

The adaptation layer uses
\begin{align}
f_{\mathrm{ref}} &= 0.12~\si{Hz},\\
\sigma_{\mathrm{ref}} &= 0.95,\\
\tau_{\mathrm{in}} &= 4.5~\si{s},\\
\tau_{\mathrm{param}} &= 7.5~\si{s},
\end{align}
together with the conservative exponent set shown directly in the configuration
code.

The simulation-backed Mahony tuning uses
\begin{align}
(2K_p)_{\mathrm{base}} &= 0.45,\\
(2K_i)_{\mathrm{base}} &= 0.015,\\
(2K_p)_{\mathrm{calm}} &= 0.90,\\
(2K_p)_{\mathrm{rough}} &= 0.35,\\
(2K_i)_{\mathrm{calm}} &= 0.025,\\
(2K_i)_{\mathrm{rough}} &= 0.010.
\end{align}

The trust references are
\begin{align}
\sigma_{\mathrm{ref}}^{\mathrm{Mahony}} &= 0.18~\si{m/s^2},\\
\varepsilon_{\mathrm{norm,ref}} &= 0.08,\\
\varepsilon_{\mathrm{innov,ref}} &= 0.12,
\end{align}
with Mahony gain-schedule smoothing time constant
\begin{equation}
\tau_g = 2.0~\si{s},
\end{equation}
and minimum accelerometer trust
\begin{equation}
\gamma_{\min} = 0.05.
\end{equation}

\subsection{Frame Mapping in the Simulation Adapter}

The simulation runner provides body NED data. The adapter converts that to the
Mahony body convention by
\begin{equation}
(x,y,z)_{\mathrm{NED}}
\mapsto
(x,-y,-z)_{\mathrm{Mahony}}.
\end{equation}

So north stays \(+X\), east becomes \(-Y\), and down becomes \(-Z\),
making gravity \(+Z\) in Mahony’s frame.

For RMS reporting back in the simulator convention, the adapter maps Mahony
Euler outputs back to the nautical convention before comparing against
reference attitude traces.

\subsection{Quality Gates Used in the Example}

The simulation code checks the last \SI{60}{s} of each run and applies quality
gates:
\begin{itemize}[leftmargin=1.4em]
  \item Z RMS \(< 17.9\%\) of \(H_s\) for JONSWAP cases,
  \item Z RMS \(< 16.0\%\) of \(H_s\) for PM-Stokes cases,
  \item yaw RMS \(< 5.06^\circ\) when magnetometer is used.
\end{itemize}

These are not part of the observer itself, but they are useful context for the
intended tuning regime in the supplied simulation.

\begin{table}[t]
  \centering
  \caption{Representative tuning from the simulation wrapper}
  \label{tab:tuning}
  \small
  \begin{tabular}{@{}lll@{}}
    \toprule
    Parameter & Value & Role \\
    \midrule
    \(r\) & 0.150 & base repeated pole rate \\
    \(\tau_a\) & \SI{0.68}{s} & accel LPF time constant \\
    \(\tau_d\) & \SI{49}{s} & slow trend time constant \\
    \(k_b\) & \(2.5\times10^{-5}\) & bias learning gain \\
    \(\lambda_b\) & \(3.0\times10^{-3}\) & bias leak \\
    \(f_{\mathrm{ref}}\) & \SI{0.12}{Hz} & reference scheduling frequency \\
    \(\sigma_{\mathrm{ref}}\) & 0.95 & reference accel sigma \\
    \((2K_p)_{\mathrm{base}}\) & 0.45 & Mahony base proportional gain \\
    \((2K_i)_{\mathrm{base}}\) & 0.015 & Mahony base integral gain \\
    \(g\) & \SI{9.80665}{m/s^2} & gravity magnitude \\
    \bottomrule
  \end{tabular}
\end{table}

% ---------------------------------------------------------------
\section{Overall Interpretation}

The complete method is a cascaded observer for vertical motion in rough seas:
\begin{enumerate}[leftmargin=1.5em]
\item Mahony AHRS estimates orientation from gyroscope integration corrected by
      gravity and optional magnetometer aiding.
\item Body-frame specific force is rotated into a \(Z\)-up world frame and
      converted to vertical inertial acceleration by subtracting gravity.
\item A policy-selected frequency tracker provides a dominant motion-frequency
      proxy and confidence information.
\item A vertical PII observer integrates the vertical acceleration to
      displacement while using repeated-pole damping, displacement-integral
      feedback, and an optional very-slow bias channel to prevent unbounded
      drift.
\item The observer bandwidth and bias-learning rates are scheduled from the
      estimated motion frequency and acceleration sigma, so calmer conditions
      permit tighter tracking while rougher conditions slow the loop and reduce
      sensitivity to spurious bias inference.
\item The Mahony gains are also reduced when accelerometer trust is poor, so
      tilt correction does not chase wave-induced specific force as if it were
      gravity.
\end{enumerate}

This design avoids a full stochastic state-space model yet still introduces
physically meaningful structure: attitude resolves gravity, the frequency
proxy provides a representative time scale, and the adaptive PII core shapes
heave dynamics to that estimated motion regime.

\section{Results Plots}

\begin{figure}[H]\centering
  \resizebox{\textwidth}{!}{\input{w3d_nonkalman_pmstokes_medium.pgf}}
  \caption{Adaptive PII/Mahony attitude (roll, pitch, yaw) --- PM Stokes, medium $H_s$.}
  \label{fig:w3d_nonkalman_pmstokes_medium_angles}
\end{figure}

\begin{figure}[H]\centering
  \resizebox{\textwidth}{!}{\input{w3d_nonkalman_pmstokes_medium_zkin.pgf}}
  \caption{Adaptive PII/Mahony vertical kinematics (displacement, velocity, acceleration with learned envelopes) --- PM Stokes, medium $H_s$.}
  \label{fig:w3d_nonkalman_pmstokes_medium_zkin}
\end{figure}

\begin{figure}[H]\centering
  \resizebox{\textwidth}{!}{\input{w3d_nonkalman_pmstokes_medium_tuner.pgf}}
  \caption{Adaptive PII/Mahony frequency and auto-scheduling traces ($f$, accel std / $\sigma_a$, $\tau$, scheduling gain) --- PM Stokes, medium $H_s$.}
  \label{fig:w3d_nonkalman_pmstokes_medium_tuner}
\end{figure}


% ---------------------------------------------------------------
\section{Summary}

The code implements a compact, layered vertical-motion estimation stack.

First, \texttt{VerticalPIIObserver} is a repeated-real-pole PII observer with
optional very-slow bias estimation. It is intentionally conservative,
non-oscillatory, and inexpensive to run.

Second, \texttt{AdaptiveVerticalPII} adds acceleration sigma estimation, a
policy-selected frequency tracker, and runtime scheduling of the observer
parameters. External displacement-frequency scheduling is preferred, while the
internal frequency proxy provides fallback scheduling.

Third, \texttt{AdaptiveVerticalPIIMahony} uses Mahony AHRS to supply the
required vertical inertial acceleration and can also adapt Mahony gains from
accelerometer trust and sea-state roughness indicators.

Together, these layers form a practical embedded heave-estimation pipeline
that is easy to tune, robust to moderate motion variability, and compatible
with several lightweight frequency-tracker backends.

\end{document}
